\subsection{Probability}\label{sec:theory_general_probability}

Probability is your best friend when dealing with real-world
electronic components. All electronic device properties suffer
from variations. There are two types of variations:
\begin{enumerate}
  \item \textbf{Systematic Variations:} These are predictable
    variations that can be measured and treated for by designing
    the circuit accordingly. For example, offset voltages in
    operational amplifiers can be measured and compensated for
    by adding a counteracting voltage source in series.
  \item \textbf{Random Variations:} These are unpredictable
    variations that follow a statistical distribution (often
    assumed to be normal/Gaussian). It is impossible to
    eliminate them, but their impact can be minimized through
    careful design practices. There are two main scenarios for
    analyzing random variations:
    \begin{enumerate}
      \item Pre-production: Using tools like Monte Carlo
        simulations to nudge various parameters within their
        tolerance ranges and see how the overall circuit performs.
      \item Post-production: Measuring the actual values of
        components after manufacturing to understand how they
        deviate from their nominal specifications.
    \end{enumerate}
    In both scenarios, statistical analysis is performed by
    measuring the property of interest (like resistance,
    capacitance, inductance, threshold voltage, gain, etc.)
    multiple times across different samples and plotting the
    results to observe the distribution of values. There are
    a few mathematical tools available to help analyze and
    interpret these distributions (see below).
\end{enumerate}

\subsubsection{Variation Tools}
It is common to assume that Random Variations follow a Normal
(Gaussian) Distribution (although this may not always be the
case!). One can use the following tools to analyze and
interpret the data collected from measurements or simulations:
\begin{itemize}
  \item \textbf{Mean (Average):} The mean is the average value
    of all measured component values.\[M = \frac{1}{N} \sum_{i=1}^{N} x_i\]
    where \(N\) is the total number of measurements and \(x_i\)
    represents each individual measurement. This gives a
    central value around which the component values are
    distributed. A designer usually designs for the mean value.
  \item \textbf{Variance (\(Var\)):} Variance measures how
    much the component values deviate from the mean. It is
    calculated as the average of the squared differences from
    the mean.\[Var = \frac{1}{N} \sum_{i=1}^{N} (x_i - M)^2\]
    For independent \(X\), \(Y\) the variances add,
    \(Var(X \pm Y) = Var(X) + Var(Y)\): expanding
    \(M[((X - M_X) \pm (Y - M_Y))^2] = Var(X) \pm 2M[(X - M_X)(Y - M_Y)] + Var(Y)\) (\(M[\cdot]\) the mean), the
    cross term is zero for independent variables.
  \item \textbf{Standard Deviation (\(SD\)):} The
    standard deviation is the square root of the variance and
    provides a measure of how tightly or loosely the
    component values are clustered around the mean. A smaller
    standard deviation indicates that the values are closely
    clustered, while a larger standard deviation indicates
    more spread out values.\[SD = \sqrt{Var}\]
  \item \textbf{Histogram:} A histogram is a graphical
    representation of the distribution of component values. It
    shows the frequency of occurrence of different values
    within specified ranges (bins). This helps visualize how
    the component values are distributed.
  \item \textbf{Cumulative Distribution Function (CDF):} The CDF
    represents the probability that a random variable takes on
    a value less than or equal to a specific value. It is useful
    for determining the likelihood of a component falling within
    a certain range. \[F(x) = P(X \leq x)\]
  \item \textbf{Probability Density Function (PDF):} The PDF
    describes the relative likelihood of a random variable to
    take on a specific value. It is the derivative of the CDF and
    provides insights into the distribution of component values.
    \[f(x) = \frac{dF(x)}{dx}\]
\end{itemize}

\subsubsection{Variation Mitigations}
For an analog integrated circuit designer, these tricks can help
mitigate the impact of variations:
\begin{itemize}
  \item \textbf{Redundancy:} Sometimes, Adding a few extra components
    that can be laser trimmed post-production can help achieve the
    desired performance despite variations.
  \item \textbf{Common Centroid Layout:} A common centroid layout
    technique can help reduce the impact of systematic variations
    by averaging out the effects across multiple components.
    In other words, placing components in a symmetrical pattern
    can help reduce the impact of variations caused by gradients
    in temperature, doping concentration, or other factors across
    the chip.
  \item \textbf{Area Scaling:} M. Pelgrom discovered that many
    variations decrease as the area of the component increases.
    This is known as Pelgrom's Law. By increasing the size of
    components, you can reduce the impact of random variations
    (by approximately the square root of the area):
    \[
      SD(dP) = \frac{K_{dP}}{\sqrt{W L}} \hspace{0.5cm} \Rightarrow \hspace{0.5cm}
      \frac{SD(dP)\big|_{4WL}}{SD(dP)\big|_{WL}} = \frac{\frac{K_{dP}}{\sqrt{4WL}}}{\frac{K_{dP}}{\sqrt{WL}}} = \frac{\sqrt{WL}}{\sqrt{4WL}} = \frac{1}{\sqrt{4}} = \frac{1}{2}
    \]
    where \(dP\) is the mismatch of a parameter between two identical devices of area
    \(WL\) and \(K_{dP}\) its Pelgrom coefficient: four times the area halves the
    spread. The down
    side is that larger components take up more space and may
    have higher parasitic capacitances, which can affect circuit
    performance.
  \item \textbf{Dummy Structures:} Adding dummy structures around
    critical components can help create a more uniform manufacturing
    environment, reducing variations caused by edge effects.
\end{itemize}

